% Generic bilingual structure. It does not claim compliance with any university.
\documentclass[UTF8,11pt,oneside,fontset=fandol]{ctexbook}
\usepackage[a4paper,margin=2.7cm]{geometry}
\usepackage{xcolor}
\usepackage{booktabs}
\usepackage{tabularx}
\usepackage{fancyhdr}
\usepackage[hidelinks,bookmarksnumbered]{hyperref}

\definecolor{accent}{HTML}{8A1538}
\pagestyle{fancy}
\fancyhf{}
\fancyhead[L]{\small 双语学位论文 / Bilingual Thesis}
\fancyhead[R]{\small\thepage}
\renewcommand{\headrulewidth}{0.4pt}
\ctexset{
  chapter = {
    format = \Huge\bfseries\color{accent},
    name = {第,章},
    number = \chinese{chapter}
  },
  section = {format = \Large\bfseries\color{accent}}
}
\setlength{\parindent}{2em}

\begin{document}

\begin{titlepage}
\centering
\vspace*{0.5cm}
{\Large 学位论文 / THESIS}\\[1.2cm]
{\color{accent}\rule{0.86\textwidth}{1.5pt}}\\[1cm]
{\Huge\bfseries 面向科学发现的\\
可信机器学习方法\par}
\vspace{0.6cm}
{\LARGE\bfseries Trustworthy Machine Learning\\
for Scientific Discovery\par}
\vspace{1.3cm}

\begin{tabular}{rl}
作者 / Author: & 林晓雨 / Xiaoyu Lin \\
学位 / Degree: & 博士 / Doctor of Philosophy \\
专业 / Program: & 计算机科学 / Computer Science \\
导师 / Supervisor: & 陈明教授 / Professor Ming Chen \\
\end{tabular}

\vfill
{\Large 示例大学 / Example University}\\[0.3cm]
2026年7月 / July 2026
\end{titlepage}

\frontmatter
\chapter*{使用说明 / Compliance note}
本模板提供通用的中英文论文结构，不代表任何学校的官方格式。提交前，请逐项核对学校最新的页边距、字体、封面、声明、盲审及电子归档要求。

This template provides a generic Chinese-English structure and is not an
official format for any institution. Before submission, verify the current
rules for margins, fonts, title pages, declarations, blind review, and
electronic archiving.

\chapter*{中文摘要}
科学机器学习模型常在训练分布内表现良好，却在实验设备、采样策略或研究地点变化时失效。
本文提出一套面向科学数据的可信学习框架，将分布偏移检测、不确定性校准和可复现实验记录统一
到同一工作流程中。在三个跨领域数据集上的实验表明，该框架能够在保持预测精度的同时显著
降低错误置信度。

\noindent\textbf{关键词：}可信机器学习；科学计算；不确定性；可复现研究

\chapter*{Abstract}
Scientific machine-learning models often perform well in distribution but fail
when instruments, sampling policies, or study sites change. This thesis unifies
shift detection, uncertainty calibration, and reproducible experiment records
in one workflow. Across three cross-domain datasets, the framework reduces
confident errors while preserving predictive accuracy.

\noindent\textbf{Keywords:} trustworthy machine learning; scientific computing;
uncertainty; reproducible research

\tableofcontents

\mainmatter
\chapter{绪论 / Introduction}
科学数据的采集过程会改变模型所见的分布。因此，可靠的模型不仅需要预测结果，还需要说明
适用范围、数据来源和不确定性。

Scientific data collection changes the distribution observed by a model.
Reliable systems must therefore report not only predictions, but also scope,
data provenance, and uncertainty.

\section{研究问题 / Research questions}
\begin{tabularx}{\textwidth}{@{}lX@{}}
\toprule
\textbf{问题 / Question} & \textbf{验证方式 / Evaluation} \\
\midrule
分布偏移何时可被检测？ & 跨设备和跨地点的预注册外部验证 \\
When is shift detectable? & Pre-registered validation across devices and sites \\
置信度何时可信？ & 分组校准误差与覆盖率分析 \\
When is confidence reliable? & Group calibration error and coverage analysis \\
\bottomrule
\end{tabularx}

\backmatter
\chapter*{数据与代码 / Data and code availability}
填写数据集版本、许可协议、永久链接、软件环境及完整复现实验命令。受限数据应提供访问流程和
可公开运行的合成数据验证路径。

Record dataset versions, licenses, permanent identifiers, software
environments, and exact reproduction commands. For restricted data, provide an
access process and a public synthetic-data validation path.

\end{document}
